\documentclass[../../main.tex]{subfiles}

\begin{document}

\chapter{Pengantar Assembly dan Sistem Bilangan}
\label{chap:pengantar-assembly}

% =========================================================
% OUTCOME PEMBELAJARAN KHUSUS BAB
% =========================================================

\begin{subcpmkbox}
\vspace{1.1em}
\SubCPMK{1.1} --- Menjelaskan peran bahasa assembly serta mengonversi
bilangan biner, desimal, oktal, dan heksadesimal, termasuk membedakan
bilangan bertanda dan tidak bertanda.

\vspace{0.4em}
\textbf{CPMK yang didukung:} \CPMK{01} \quad\textbullet\quad
\textbf{CPL yang didukung:} \CPL{02}
\end{subcpmkbox}

\begin{tujuanbox}
\vspace{1.1em}
Setelah mempelajari bab ini, mahasiswa mampu:
\begin{enumerate}
    \item menjelaskan kedudukan bahasa assembly di antara bahasa mesin
          dan bahasa tingkat tinggi;
    \item menjelaskan peran bahasa assembly dalam pemrograman sistem
          dan kendali perangkat keras;
    \item mengonversi bilangan biner, desimal, oktal, dan heksadesimal
          beserta langkah kerjanya;
    \item menuliskan literal bilangan yang diterima oleh assembler;
    \item membedakan bilangan bertanda dan tidak bertanda serta
          menerapkan \eng{two's complement};
    \item menjelaskan hubungan kode ASCII dengan data yang disimpan di
          memori.
\end{enumerate}
\end{tujuanbox}

\begin{materibox}
\vspace{1.1em}
\begin{enumerate}
    \item Tingkat abstraksi bahasa pemrograman.
    \item Peran dan kedudukan bahasa assembly 8086.
    \item Sistem bilangan biner, desimal, oktal, dan heksadesimal.
    \item Konversi antarjenis bilangan.
    \item Penulisan literal bilangan pada assembler.
    \item Bilangan bertanda, bilangan tidak bertanda, dan
          \eng{two's complement}.
    \item Kode ASCII dan representasi data satu byte.
\end{enumerate}
\end{materibox}

\subfile{section-01}
\subfile{section-02}
\subfile{section-03}

\section{Contoh Program}
\label{sec:contoh-bab3}

Bagian ini menyajikan tiga contoh program yang dirakit dan dijalankan
menggunakan Emu8086 atau TASM. Ketiganya sengaja dibuat sesingkat
mungkin agar seluruh baris dapat dipahami satu per satu.

\subsection{Contoh 1: Mencetak Teks dengan Layanan DOS}

Program pada Kode~\ref{lst:hello} adalah program assembly pertama yang
biasa dipakai: mencetak satu baris teks ke layar.

\lstinputlisting[
    language=MASM,
    caption={Program \file{hello.asm} mencetak teks ke layar.},
    label={lst:hello}
]{\CodeDir/hello.asm}

Penjelasan setiap bagian program:

\begin{itemize}
    \item \lstasm|.MODEL SMALL| memberi tahu assembler bahwa program
          memakai model memori kecil.
    \item \lstasm|.DATA| membuka segmen data; di dalamnya terdapat
          variabel \code{pesan} berupa deretan byte yang diakhiri
          \code{0Dh}, \code{0Ah} (pindah baris), dan \code{'\$'}.
    \item \lstasm|panjang EQU \$ - pesan| menghitung panjang string
          secara otomatis dari selisih posisi sekarang dengan awal
          label \code{pesan}.
    \item \lstasm|MOV AH, 09h| dan \lstasm|LEA DX, pesan| menyiapkan
          argumen layanan DOS: \reg{AH} berisi nomor fungsi, \reg{DX}
          berisi alamat string.
    \item \lstasm|INT 21h| menyerahkan kendali ke layanan DOS untuk
          mencetak string.
    \item \lstasm|MOV AH, 4Ch| dan \lstasm|INT 21h| mengakhiri program.
\end{itemize}

\begin{session}[caption={Keluaran program \file{hello.asm}.}, label={lst:keluaran-hello}]
Halo, dunia!
\end{session}

\subsection{Contoh 2: Program \file{.COM} dengan \dir{ORG 100h}}

Bentuk lain yang lebih ringkas adalah program \file{.COM}. Karena
seluruh kode, data, dan stack berada dalam satu segmen, inisialisasi
\reg{DS} tidak diperlukan. Program ditunjukkan pada Kode~\ref{lst:com-org}.

\begin{lstlisting}[
    language=MASM,
    caption={Program \file{.COM} ringkas dengan \dir{ORG 100h}.},
    label={lst:com-org}
]
.MODEL TINY
.CODE
ORG 100h            ; program .COM dimulai pada offset 100h

START:
    MOV AH, 09h     ; cetak string
    MOV DX, OFFSET pesan
    INT 21h

    MOV AH, 4Ch     ; keluar
    MOV AL, 0
    INT 21h

pesan DB 'Selamat datang di assembly!', 0Dh, 0Ah, '$'
END START
\end{lstlisting}

\begin{infobox}
Perhatikan bahwa pada program \file{.COM} tidak ada
\lstasm|MOV AX, @DATA| maupun \lstasm|MOV DS, AX|. DOS sudah
menetapkan \reg{DS} ke segmen program, sehingga variabel dapat
langsung dialamatkan.
\end{infobox}

\subsection{Contoh 3: Menampilkan Karakter Satu per Satu}

Layanan DOS fungsi \hex{02} mencetak satu karakter yang berada pada
register \reg{DL}. Contoh ini menunjukkan bagaimana sebuah huruf
dapat dicetak, kemudian diubah, lalu dicetak kembali
(Kode~\ref{lst:karakter}).

\begin{lstlisting}[
    language=MASM,
    caption={Mencetak karakter dengan \code{INT 21h} fungsi \hex{02}.},
    label={lst:karakter}
]
.MODEL SMALL
.STACK 100h

.CODE
START:
    MOV AX, @DATA
    MOV DS, AX

    MOV AH, 02h     ; fungsi 02h : cetak karakter pada DL
    MOV DL, 'A'
    INT 21h

    MOV DL, 'B'
    INT 21h

    MOV DL, 'C'
    INT 21h

    MOV AH, 4Ch
    MOV AL, 0
    INT 21h

END START
\end{lstlisting}

\begin{session}[caption={Keluaran program pada Kode~\ref{lst:karakter}.}]
ABC
\end{session}

Ketiga contoh di atas menunjukkan pola yang sama: siapkan argumen
pada register, panggil layanan melalui \code{INT 21h}, lalu akhiri
program. Pola inilah yang akan dipakai ulang pada bab-bab berikutnya.

\section{Praktikum}
\label{sec:praktikum-bab3}

\begin{praktikumbox}
\vspace{1.1em}
\textbf{Judul:} Konversi bilangan dan pemeriksaan data di memori.\\[0.3em]
\textbf{Sub-CPMK:} 1.1 \quad\textbullet\quad
\textbf{Alokasi waktu:} 100 menit \quad\textbullet\quad
\textbf{Bentuk:} individu
\end{praktikumbox}

\subsection{Tujuan Praktikum}

Setelah menyelesaikan praktikum ini, mahasiswa mampu mengonversi
bilangan antar basis, menuliskan nilai dalam beberapa basis pada
program assembly, serta memeriksa dan menafsirkan isi memori dengan
DEBUG dan emu8086.

\subsection{Alat dan Bahan}

\begin{itemize}
    \item Komputer dengan sistem operasi Windows.
    \item emu8086.
    \item DOSBox dengan TASM, TLINK, dan DEBUG.
    \item Berkas \file{bilangan.asm} dari \file{code/asm/}.
\end{itemize}

\subsection{Langkah Kerja}

\begin{enumerate}
    \item \textbf{Konversi manual.}
    \begin{enumerate}
        \item Salin tabel berikut ke lembar kerja.
        \item Lengkapi kolom biner, oktal, dan heksadesimal tanpa
              memakai kalkulator.
    \end{enumerate}

    \begin{center}
    \begin{tabular}{@{}cccc@{}}
        \toprule
        \textbf{Desimal} & \textbf{Biner} & \textbf{Oktal} &
        \textbf{Heksadesimal} \\
        \midrule
        37  & & & \\
        45  & & & \\
        100 & & & \\
        156 & & & \\
        200 & & & \\
        243 & & & \\
        255 & & & \\
        \bottomrule
    \end{tabular}
    \end{center}

    \item \textbf{Menulis program data.}
    \begin{enumerate}
        \item Buka \file{bilangan.asm} pada editor emu8086 atau editor
              teks.
        \item Perhatikan empat baris data: \code{desimal},
              \code{biner}, \code{heksa}, dan \code{karakter}.
        \item Jelaskan pada lembar kerja mengapa tiga baris pertama
              menghasilkan byte yang sama.
    \end{enumerate}

    \item \textbf{Perakitan dan pengaitan.}
    \begin{enumerate}
        \item Buka DOSBox, lalu masuk ke direktori kerja.
        \item Rakit dengan perintah \file{tasm bilangan.asm}.
        \item Link dengan perintah \file{tlink bilangan.obj}.
        \item Jalankan program dengan perintah \file{bilangan}.
    \end{enumerate}

    \item \textbf{Pemeriksaan memori dengan DEBUG.}
    \begin{enumerate}
        \item Jalankan \file{debug bilangan.exe}.
        \item Ketik \file{d ds:0} untuk menampilkan isi segmen data.
        \item Catat delapan byte pertama yang muncul.
        \item Tandai byte yang mewakili \code{desimal}, \code{biner},
              \code{heksa}, dan \code{karakter}.
    \end{enumerate}

    \item \textbf{Modifikasi data.}
    \begin{enumerate}
        \item Ubah nilai data menjadi:
              \code{desimal DB 37}, \code{biner DB 11011011b},
              \code{heksa DB 0DBh}, dan \code{karakter DB 'Z'}.
        \item Rakit ulang, lalu periksa kembali isi memori.
        \item Tafsirkan byte \hex{DB} sebagai bilangan tidak bertanda
              dan sebagai bilangan bertanda 8 bit.
    \end{enumerate}

    \item \textbf{Pengamatan dengan emu8086.}
    \begin{enumerate}
        \item Buka program yang sama pada emu8086.
        \item Jalankan dengan \file{Single Step} sampai instruksi
              \lstasm|MOV DS, AX| selesai.
        \item Buka jendela \file{Memory} pada alamat \code{DS:0000}.
        \item Bandingkan tampilan emu8086 dengan hasil DEBUG.
    \end{enumerate}
\end{enumerate}

\subsection{Hasil yang Dikumpulkan}

\begin{enumerate}
    \item Tabel konversi bilangan yang telah dilengkapi.
    \item Berkas \file{bilangan.asm} versi modifikasi.
    \item Tangkapan layar hasil \file{d ds:0} pada DEBUG.
    \item Tangkapan layar jendela memori emu8086.
    \item Analisis singkat: mengapa \hex{F3} dapat berarti 243 atau
          $-13$, dan mengapa \code{'A'} disimpan sebagai \hex{41}.
\end{enumerate}

\begin{asesmenbox}
\vspace{1.1em}
\textbf{Asesmen praktikum 1 (bobot 5\% dari komponen praktikum).}

\begin{tabular}{@{}p{6.2cm}cp{4.6cm}@{}}
    \toprule
    \textbf{Indikator} & \textbf{Bobot} & \textbf{Kriteria} \\
    \midrule
    Ketepatan konversi bilangan & 35\% &
    Seluruh baris tabel benar atau kesalahan dapat dijelaskan \\
    \addlinespace
    Program data berhasil dirakit dan dijalankan & 25\% &
    Berkas \file{.EXE} terbentuk dan program selesai tanpa \eng{hang} \\
    \addlinespace
    Ketepatan pembacaan isi memori & 25\% &
    Empat byte data dikenali dan ditafsirkan dengan benar \\
    \addlinespace
    Analisis bilangan bertanda dan ASCII & 15\% &
    Penjelasan \hex{F3} dan \code{'A'} tepat serta beralasan \\
    \bottomrule
\end{tabular}
\end{asesmenbox}

\section{Latihan}
\label{sec:latihan-bab3}

\begin{enumerate}
    \item \textbf{Konseptual.} Jelaskan perbedaan bahasa mesin, bahasa
          assembly, dan bahasa tingkat tinggi. Sertakan satu contoh
          operasi sederhana yang ditulis pada ketiga tingkat tersebut.

    \item \textbf{Konseptual.} Mengapa hubungan antara bahasa assembly
          dan bahasa mesin disebut \eng{one-to-one}? Bandingkan dengan
          hubungan bahasa tingkat tinggi dan bahasa mesin.

    \item \textbf{Konversi.} Ubah 156 desimal menjadi biner, oktal,
          dan heksadesimal. Tuliskan langkah pembagian atau
          pengelompokan bit yang Anda pakai.

    \item \textbf{Konversi dan tafsir.} Diberikan \bin{10110110}.
          Tentukan nilai desimalnya sebagai bilangan tidak bertanda
          dan sebagai bilangan bertanda 8 bit. Jelaskan perbedaan
          hasilnya.

    \item \textbf{Two's complement.} Tentukan bentuk 8 bit dari $-45$
          dengan langkah membalik bit dan menambahkan 1. Periksa
          hasilnya dengan mengurangkan 256 dari nilai tidak
          bertandanya.

    \item \textbf{Literal assembler.} Jelaskan mengapa \code{0F3H}
          ditulis dengan nol di depan, sedangkan \code{F3H} dapat
          menimbulkan kesalahan. Sebutkan pula dua kesalahan lain
          dalam penulisan literal bilangan.

    \item \textbf{Penulisan data.} Tulislah potongan segmen data yang
          mendefinisikan nilai 156 dalam bentuk desimal, biner, dan
          heksadesimal, lalu satu byte karakter \code{'Z'}.
          Prediksikan empat byte yang akan tampil pada DEBUG.

    \item \textbf{Analisis.} Sebuah pemeriksaan memori menghasilkan
          \code{C8 C8 C8 5A}. Jelaskan makna setiap byte, baik sebagai
          bilangan tidak bertanda maupun sebagai karakter.
\end{enumerate}

\begin{refleksibox}
\vspace{1.1em}
Jawab pertanyaan berikut pada jurnal praktikum Anda.
\begin{enumerate}
    \item Bagian konversi mana yang paling sulit Anda kerjakan?
          Langkah apa yang membantu Anda menyelesaikannya?
    \item Setelah memeriksa memori dengan DEBUG, apakah pemahaman Anda
          tentang hubungan bilangan dan data berubah? Jelaskan.
    \item Konsep apa yang masih perlu Anda ulangi sebelum mempelajari
          arsitektur 8086 pada bab berikutnya?
\end{enumerate}
\end{refleksibox}

\section{Rangkuman}
\label{sec:rangkuman-bab3}

Pada bab ini telah dibahas:

\begin{itemize}
    \item Bahasa pemrograman dibedakan menurut tingkat abstraksinya:
          bahasa mesin, bahasa assembly, dan bahasa tingkat tinggi.
          Bahasa assembly berada tepat di atas bahasa mesin dan
          bersifat \eng{one-to-one} terhadapnya.

    \item Bahasa assembly tetap dipakai ketika dibutuhkan kendali
          langsung atas perangkat keras, ukuran program yang kecil,
          serta pemahaman cara kerja prosesor pada level instruksi.

    \item Bahasa assembly 8086 mengenal bilangan biner, oktal,
          desimal, dan heksadesimal. Setiap basis mempunyai aturan
          digit dan bobot posisi yang berbeda.

    \item Konversi dapat dilakukan dengan menjumlahkan bobot posisi,
          membagi berulang dengan basis tujuan, atau mengelompokkan
          bit menjadi tiga bit untuk oktal dan empat bit untuk
          heksadesimal.

    \item Assembler menerima literal desimal tanpa akhiran atau
          dengan akhiran \code{D}, biner dengan akhiran \code{B}, dan
          heksadesimal dengan akhiran \code{H}. Heksadesimal yang
          diawali huruf A--F harus diberi nol di depan.

    \item Bilangan tidak bertanda memakai seluruh bit sebagai
          besaran. Bilangan bertanda memakai bit paling kiri sebagai
          bit tanda dan sisanya mengikuti aturan
          \eng{two's complement}.

    \item Pola bit yang sama dapat ditafsirkan berbeda. \hex{F3}
          berarti 243 sebagai bilangan tidak bertanda dan $-13$
          sebagai bilangan bertanda 8 bit.

    \item Karakter disimpan sebagai kode ASCII satu byte. \code{'A'}
          disimpan sebagai \hex{41}, sehingga data karakter dan data
          bilangan dapat diperiksa dengan cara yang sama di memori.
\end{itemize}

\begin{checklistbox}
\vspace{1.1em}
Tandai setiap butir yang sudah Anda kuasai.

\begin{itemize}
    \item[\checkbox] Saya dapat menjelaskan tiga tingkat bahasa
          pemrograman dan kedudukan assembly.
    \item[\checkbox] Saya dapat mengonversi bilangan biner, oktal,
          desimal, dan heksadesimal.
    \item[\checkbox] Saya dapat menuliskan literal bilangan yang
          diterima assembler, termasuk nol di depan heksadesimal.
    \item[\checkbox] Saya dapat membedakan bilangan bertanda dan
          tidak bertanda.
    \item[\checkbox] Saya dapat membentuk \eng{two's complement} dan
          menafsirkan pola bit seperti \hex{F3}.
    \item[\checkbox] Saya dapat menjelaskan hubungan kode ASCII dengan
          isi memori.
    \item[\checkbox] Saya dapat membaca hasil \file{d ds:0} pada DEBUG
          dan menjelaskan setiap byte.
\end{itemize}

\vspace{0.4em}
Bila masih ada butir yang belum bertanda, ulangi bagian yang
bersesuaian sebelum melanjutkan ke bab berikutnya.
\end{checklistbox}

\begin{infobox}
\vspace{1.1em}
\textbf{Rujukan bab.} Materi pada bab ini diringkas dari
\cite{gunawan2019}, Bagian 1, hlm.~10--11 dan Bagian 14.1,
hlm.~127--128, serta \cite{sto2001}, BAB I, hlm.~6--8, BAB V,
hlm.~23--24, dan BAB XV, hlm.~78.
\end{infobox}

\section{Evaluasi}
\label{sec:evaluasi-bab3}

\subsection{Soal Pilihan Ganda}

\begin{enumerate}
    \item Bilangan biner adalah bilangan yang berbasis \dots
          \begin{enumerate}
              \item 2 dengan digit 0 dan 1
              \item 8 dengan digit 0 sampai 7
              \item 10 dengan digit 0 sampai 9
              \item 16 dengan digit 0 sampai F
          \end{enumerate}

    \item Nilai desimal dari \bin{1010} adalah \dots
          \begin{enumerate}
              \item 8 \item 10 \item 12 \item 16
          \end{enumerate}

    \item Nilai desimal dari \hex{F3} adalah \dots
          \begin{enumerate}
              \item 153 \item 243 \item 253 \item 263
          \end{enumerate}

    \item Pola 8 bit \bin{11111111} sebagai bilangan bertanda
          bernilai \dots
          \begin{enumerate}
              \item 255 \item 127 \item $-1$ \item $-128$
          \end{enumerate}

    \item Penulisan heksadesimal yang benar pada assembler adalah
          \dots
          \begin{enumerate}
              \item \code{A12H} \item \code{0A12H}
              \item \code{A12} \item \code{0A12D}
          \end{enumerate}
\end{enumerate}

\subsection{Soal Uraian}

\begin{enumerate}
    \item Ubah 200 desimal menjadi biner, oktal, dan heksadesimal.
          Tuliskan langkah pembagian dan pengelompokan bit yang Anda
          gunakan.

    \item Jelaskan perbedaan bilangan bertanda dan tidak bertanda pada
          8 bit. Tentukan bentuk \eng{two's complement} dari $-37$ dan
          periksa hasilnya dengan menghitung nilai tidak bertandanya.

    \item Pola \hex{F3} dapat berarti 243 atau $-13$. Jelaskan mengapa
          hal itu mungkin, dan sebutkan peran \eng{Sign Flag} serta
          jenis instruksi dalam menentukan tafsirnya.

    \item Tulislah potongan segmen data yang mendefinisikan 156 dalam
          bentuk desimal, biner, dan heksadesimal, ditambah karakter
          \code{'Z'}. Prediksikan isi memori yang tampil pada DEBUG,
          lalu jelaskan setiap byte.
\end{enumerate}

\subsection{Kisi-kisi Penilaian}

\begin{table}[H]
    \centering
    \caption{Kisi-kisi evaluasi Bab~\ref{chap:pengantar-assembly}.}
    \label{tab:kisi-bab3}
    \small
    \begin{tabular}{@{}p{1.5cm}p{5.0cm}p{2.6cm}c@{}}
        \toprule
        \textbf{Sub-CPMK} & \textbf{Indikator yang diukur} &
        \textbf{Bentuk soal} & \textbf{Bobot} \\
        \midrule
        1.1 & Menjelaskan tingkat bahasa dan peran assembly &
        Pilihan ganda 1 & 15\% \\
        \addlinespace
        1.1 & Mengonversi bilangan antar basis &
        Pilihan ganda 2, 3; Uraian 1 & 25\% \\
        \addlinespace
        1.1 & Membedakan bilangan bertanda dan tidak bertanda &
        Pilihan ganda 4; Uraian 2, 3 & 30\% \\
        \addlinespace
        1.1 & Menulis literal dan menafsirkan isi memori &
        Pilihan ganda 5; Uraian 4 & 30\% \\
        \bottomrule
    \end{tabular}
\end{table}

\begin{checklistbox}
\vspace{1.1em}
\textbf{Checklist Pencapaian Kompetensi
Bab~\ref{chap:pengantar-assembly}}

Seorang mahasiswa dinyatakan kompeten pada Sub-CPMK 1.1 apabila:
\begin{itemize}
    \item[\checked] dapat menjelaskan kedudukan bahasa assembly di
          antara bahasa mesin dan bahasa tingkat tinggi;
    \item[\checked] dapat mengonversi bilangan biner, oktal, desimal,
          dan heksadesimal dengan langkah yang benar;
    \item[\checked] dapat menuliskan literal bilangan yang diterima
          assembler, termasuk penggunaan nol di depan;
    \item[\checked] dapat membentuk dan menafsirkan
          \eng{two's complement};
    \item[\checked] dapat menjelaskan hubungan kode ASCII dengan data
          yang disimpan di memori;
    \item[\checked] dapat membaca dan menjelaskan hasil pemeriksaan
          memori dengan DEBUG.
\end{itemize}
\end{checklistbox}


\end{document}
